\documentclass[12pt,a4paper,oneside]{article}
\usepackage[brazil]{babel}
\usepackage[utf8]{inputenc}
\usepackage{olymp}


\newlength{\exmpwidouf}
\exmpwidinf=10cm
\exmpwidouf=6cm


\setcounter{problem}{2}

\begin{document}

\begin{problem}{ccTLD}{codigo}{1}

O domínio de topo de código de país (ccTLD, do inglês \textit{country code top-level domain}), também chamado de domínio nacional de nível superior, é o domínio de topo na Internet geralmente usado ou reservado para um país. Os identificadores de ccTLD são de duas letras, como o \texttt{.br} usado no Brasil. Assim como no caso do Brasil, muitos países usam as duas letras inicias do seu nome: \texttt{.it} na Itália e \texttt{.de} na Alemanha (em alemão, \textit{Deutschland}). Outros usam a letra inicial e alguma outra letra do nome, seja por preferência ou por não estar mais disponível o código com as duas iniciais: México usa \texttt{.mx} e Paraguai usa \texttt{.py} (de \textit{Paraguay}) (enquanto Panamá usa \texttt{.pa}). E alguns são forçados a criar alternativas: Benin usa \texttt{.bj} pois \texttt{.be}, \texttt{.bn} e \texttt{.bi} já eram usados.

Nesta questão, seu programa deve verificar se um código usa as duas letras iniciais do nome do país ou a letra inicial e alguma outra letra do nome.

% \Notes

% Você pode criar uma função para o cálculo do fatorial.

\InputFile

A entrada contém vários casos de teste, cada um deles descrito em uma linha contendo um código (ponto e duas letras minúsculas) seguido de um espaço e o nome do país. O nome do país contém apenas letras e espaços, sem acentos ou cedilhas. O código sempre tem duas letras e o nome do país no máximo 50. O código \texttt{.zz} indica o fim da entrada.

\OutputFile

Para cada caso de teste da entrada, escreva: \texttt{Ok}, se o código usa as duas letras iniciais ou a letra inicial e alguma outra letra do nome informado; ou \texttt{No}, caso contrário.

% \Notes
% Preste atenção aos tipos das variáveis, pois $F(90) \approx 2^{61}$.

\Example

\begin{example}
\exmp{
.br Brasil
.de Alemanha
.de Germany
.de Deutschland
.za Africa do Sul
.za South Africa
.za Zuid Afrika
.tv Tuvalu
.cu Cuba
.py Paraguai
.py Paraguay
.ae Emirados Arabes Unidos
.ae United Arab Emirates
.zz
}{
Ok
No
No
Ok
No
No
Ok
Ok
Ok
No
Ok
No
No
}%
\end{example}

% \begin{example}
% \exmp{
% .br Brasil
% 2
% Brasil
% Brazil
% }{
% OK
% }%
% \end{example}

% \begin{example}
% \exmp{
% .de
% 3
% Alemanha
% Germany
% Deutschland
% }{
% OK
% }%
% \end{example}

% \begin{example}
% \exmp{
% .tv
% 1
% Tuvalu
% }{
% OK
% }%
% \end{example}

% \begin{example}
% \exmp{
% .tv
% 1
% Tuvalu
% }{
% OK
% }%
% \end{example}

% \begin{example}
% \exmp{
% .za
% 3
% Africa do Sul
% South Africa
% Zuid-Afrika
% }{
% OK
% }%
% \end{example}

% \textit{Problema baseado na questão ``A idade de dona Mônica'' da OBI2019 fase local.}

\end{problem}

\end{document}
